\subtituloA{Arquitecturas fundamentales de la web}

|

La web surgió como una sistema de documentos interconectados mediante hipervínculos, donde cada documento estaba identificado por una URL y residía en un servidor permanente. Esta arquitectura simple sentó las bases para la evolución posterior.

 |
   |
  \lista{
  **Servidor:** Equipo que almacena y sirve los recursos solicitados por los clientes a través de la red. |
  **Cliente (Navegador):** Interfaz que solicita recursos al servidor, los interpreta y los presenta al usuario a través de una interfaz gráfica.
} |
  \nota{Para el examen}{
¿Cuál es la diferencia fundamental entre un servidor estático y uno dinámico?

\textbf{Respuesta:} .
}


modelo-cliente-servidor"
\lista{}
